\chapter{实验结果讨论与实践建议}

本章在第四章实证结果的基础上进行讨论，重点解释显著变量和非显著变量的含义，并结合典型案例说明不同内容组合如何影响互动表现。在此基础上，进一步提出面向地方文旅账号的内容生产和运营建议。

\section{主要实验结果总结}

本文围绕 AI 文旅短视频营销效果展开实证分析，核心目标是回答“每一个指标对结果有无影响，以及具体影响多少”。描述统计、单因素差异检验、多元回归、高互动二分类 Logit、稳健性检验和典型案例分析共同表明，样本传播效果呈现明显长尾分布。722 条样本的总互动量均值为 11108.14，中位数为 907.5。少数爆款视频贡献了大量互动，多数视频互动水平较低。这种分布特征符合短视频平台的流量分配逻辑，也提示本文不能只依靠个别爆款案例判断传播规律。

单因素检验显示，AI 技术质感、视频风格、内容导向、视频主题、AI 介入程度、视频主体和 AI 介入形式在不同类别之间均存在显著差异，表明 AI 特征和内容特征在表面上都与互动表现有关。多元回归进一步揭示了这些变量的净影响：在控制其他因素后，真正稳定显著的变量主要是 AI 技术质感、视频风格、视频时长和发布距采集日天数。AI 技术质感中，普通质感相对于粗糙质感的影响约为 +223.0\%，精良质感相对于粗糙质感的影响约为 +1641.4\%；视频风格中，抽象/超现实风格相对于纪实/写实风格的影响约为 +368.1\%；log 视频时长影响约为 -37.9\%；log 发布距采集日天数影响约为 +28.8\%。

稳健性检验进一步支持了这一判断。替换点赞、评论、收藏和转发作为因变量后，AI 技术质感和抽象/超现实风格仍在多个模型中表现显著；剔除极端值后，核心变量依然稳定。高互动二分类 Logit 的补充结果也显示，AI 技术质感、抽象/超现实风格和视频时长能够较好区分高互动样本与普通样本。这说明本文的主要发现并非仅由某一种互动行为或少数极端爆款视频驱动。

如果把这些结果放回地方文旅账号的真实运营场景中理解，可以发现本文识别出的规律具有较强现实解释力。对于地方文旅账号而言，内容生产往往面临三重约束：一是更新频率高，二是预算和人力有限，三是必须在平台算法环境中争夺注意力。在这种情况下，“什么样的 AI 内容更容易获得互动”并不是抽象的学术问题，而是直接关系到账号是否值得继续投入某种内容生产方式。本文结果表明，真正值得重点关注的并不是 AI 是否被高频使用，而是 AI 是否最终转化为用户能直接感知的内容质量和风格优势。这一点使得研究结论具备了从统计结果走向运营判断的现实基础。

\section{核心影响因素讨论}

AI 技术质感是本文最重要的显著变量。结果显示，精良质感视频的传播效果远高于粗糙质感视频，普通质感视频也显著优于粗糙质感视频。用户反应并不是由“是否使用 AI”直接决定，而是由最终呈现出来的内容质量决定。

AIGC 视频质量研究强调，画面协调性、视觉完整度以及文本和画面之间的一致性会影响用户对生成内容的理解与评价\cite{lit03}；AIGC 营销研究也指出，生成内容能否形成价值感和情感连接，会进一步影响用户参与\cite{lit18}。在文旅场景中，AI 技术质感还会影响用户对目的地形象的判断。精良画面能够提升目的地吸引力，粗糙画面则可能削弱真实感和专业性。

因此，AI 技术质感显著说明，相关营销应从“技术展示逻辑”转向“内容质量逻辑”。地方文旅账号不能只强调自己使用了 AI，而要关注生成内容是否清晰、自然、具有审美价值，并能与地方文化和旅游资源形成有效连接。

这一发现还意味着，AI 对文旅短视频传播效果的作用取决于其是否真正提升了内容呈现质量。当内容策划、地方符号选择和视觉表达具备较好基础时，AI 可以增强画面表现力和内容吸引力；但如果前期选题混乱、地方元素薄弱或表达逻辑平直，单纯引入 AI 并不能自动改善传播效果。因此，技术质感显著并不只是说明画面精致度重要，也说明内容生产链条中的前端策划、视觉把控和地方意象提炼仍然需要被重视。

视频风格中，抽象/超现实风格显著正向影响传播效果。AI 技术的优势并不只是复制现实，而是创造现实拍摄难以实现的场景。传统文旅短视频常依赖自然景观、城市街景和文化场所实拍，而 AI 可以将地方符号转化为想象化、奇观化和情绪化的视觉表达。

短视频营销研究强调视觉内容和叙事吸引力对旅游意图的重要作用\cite{lit06}。在本文样本中，抽象/超现实风格能够强化视觉冲击和新奇感，更容易触发用户“没见过”“想转发”“想评论”的互动动机。与纪实/写实风格相比，超现实风格更符合短视频平台中快速吸引注意力的传播逻辑。

但这并不意味着所有文旅视频都应追求夸张化表达。抽象/超现实风格仍需建立在地方文化符号和旅游资源基础上，否则容易变成与目的地无关的视觉炫技。有效的 AI 风格创新应当是在真实地方特色基础上的想象力延展。

从用户互动表现看，抽象/超现实风格的优势还在于更容易形成评论和转发话题。相较于普通景观展示，这类视频更容易引发“这是哪里”“这个画面怎么做出来的”“原来某地还能这样表达”等讨论，从而带动评论和转发。对于以公共传播为目标的文旅账号而言，这一点较为重要，因为内容不仅需要被用户看到，还需要促使用户产生进一步评论、收藏或转发行为。当然，如果超现实表达与地方文化脱节，用户也可能只记住特效而忘记目的地本身。因此，风格创新的关键并不是一味夸张，而是让想象性表达始终围绕地方识别度展开。

视频时长显著负向影响传播效果，说明过长视频会降低互动表现。短视频平台用户通常处于快速浏览状态，内容需要在较短时间内建立吸引力。如果视频铺垫过长、节奏拖沓，即使使用 AI 技术，也可能无法转化为点赞、评论或转发。

典型案例也呈现出相同特征。浙江“地标建筑热到撒腿就跑”视频时长仅 22 秒，却在开头快速呈现地标拟人化和高温热点之间的冲突，视觉符号明确，节奏直接，容易形成记忆点。此类内容需要把技术亮点、地方符号和情绪表达集中呈现，而不是简单延长视频时长。

发布距采集日天数显著正向影响传播效果，主要反映互动积累效应。视频从发布到采集之间经历的时间越长，越有机会被推荐、观看和互动，因此总互动量更高。这一变量不是内容质量变量，但必须在模型中控制，否则较早发布视频的互动积累可能会被误判为内容特征优势。

\section{非显著变量讨论}

主回归中，AI 介入形式、AI 介入程度、视频主体、视频主题和内容导向未呈现显著影响。这一结果需要结合单因素检验和回归模型共同理解。

AI 介入形式不显著，表明 AI 用在声音、画面还是混合层面，本身并不是决定传播效果的关键。AI 画面和混合形式在表面上可能更容易吸引注意，但如果画面质感一般、风格缺乏新意，仍然难以获得高互动。用户并不会因为某个视频属于“AI 画面”或“混合 AI”就自动互动。

AI 介入程度不显著，说明深度 AI 介入并不必然优于浅度或中度介入。深度 AI 介入可能带来更强的新奇感，也可能带来更明显的生成痕迹和真实感问题。AI 内容可信度研究表明，用户对 AI 内容的接受受到可信度和感知质量影响\cite{lit14}。因此，AI 介入程度更像是生产过程变量，而不是直接决定用户互动的结果变量。

视频主体和视频主题不显著，意味着在控制其他变量后，“拍人、拍景还是拍物”“自然风光还是历史人文”等主题分类本身不足以解释传播差异。相同主题下，视频可能因为技术质感、风格和节奏不同而表现出完全不同的互动结果。文旅短视频营销不能只依赖资源禀赋，还需要依赖内容表达。

内容导向在单因素检验中显著，但回归中不显著，表明娱乐审美型内容的优势可能已被视频风格和技术质感吸收。也就是说，娱乐审美型视频之所以更容易获得高互动，可能是因为它们更常采用精良 AI 画面、抽象/超现实风格和更短节奏，而不是“娱乐审美型”标签本身直接起作用。

这些非显著结果同样具有研究价值。本文设置较多变量，是为了完整控制 AI 使用方式、内容主题、视觉表达和发布时间条件，而不是预设每个变量都应在回归中显著。单因素检验说明许多分类变量在表面上存在互动差异，多元回归则进一步回答在同时控制其他因素后，哪些变量仍具有独立解释力。因此，部分变量不显著并不意味着它们没有内容意义，而是说明其影响可能已经被 AI 技术质感、视频风格、视频时长和发布时间等更直接的因素吸收。

换言之，非显著并不等于“无用”，而更接近“不是最直接的驱动因素”。例如，视频主题和内容导向仍然会影响创作方向、用户预期和地方传播定位，只是这些影响往往需要通过更具体的呈现质量与风格表达才能真正转化为互动结果。对于文旅账号运营而言，这意味着选题当然重要，但选题的传播效果最终仍取决于它被怎样视觉化、节奏化和平台化。这个判断也有助于避免运营实践中常见的误区，即把传播不佳简单归咎于“题材选错了”，而忽视了表达方式本身的问题。

\section{稳健性结果讨论}

稳健性检验在本文中主要用于判断核心结论是否依赖某一特定结果指标、某一类极端值样本、某一特定时间累计口径或地区差异。结果显示，AI 技术质感、视频风格和视频时长等核心变量在替换因变量、调整时间口径、加入省份固定效应和处理极端值后仍保持较强稳定性，说明本文结论具有一定稳健性。

尤其需要强调的是，考虑到互动量会受到采集时点影响，本文将时间变量统一调整为发布距采集日天数，并在稳健性检验中引入日均互动量这一辅助口径。该处理表明，时间积累因素确实会影响总互动量大小，但并未改变 AI 技术质感和抽象/超现实风格是核心影响因素这一总体结论。

稳健性检验结果表明，AI 技术质感、视频风格与视频时长并非某一组特殊样本或某一种结果变量下的偶然产物，而是在不同互动指标、极端值处理方式和时间口径下相对稳定出现的因素。

\section{典型案例分析}

本文在完成回归分析与稳健性讨论之后，选取高互动和低互动样本中的典型视频进行补充分析。案例选取围绕 AI 技术质感、视频风格、视频时长和内容导向等核心变量展开，重点比较不同内容组合下的互动表现。

浙江“地标建筑热到撒腿就跑”视频总互动量达到 485737。该视频以高温天气为切入点，将浙江地标建筑拟人化处理，形成“地标也热到逃跑”的超现实表达。视频技术质感为精良，风格为抽象/超现实，时长仅 22 秒。它较好体现了精良质感、奇观化表达和短时长节奏之间的结合，说明高互动内容往往不是简单展示景观，而是把地方符号转化为可被平台用户快速理解和传播的视觉梗。

广西相关视频总互动量达到 514000，视频以“AI眼中的广西”为主题，采用超现实方式呈现地方形象。该案例说明，AI 技术能够突破现实拍摄限制，把地方景观转化为更具想象力的视觉内容，从而提升用户的新奇感和互动意愿。这与前文中“抽象/超现实风格显著正向影响互动效果”的统计结果相互印证。

天津某 AI 节气主题视频总互动量仅为 10。该视频虽然具备 AI 介入，但技术质感较为粗糙，传播效果较弱。这说明深度或混合 AI 介入并不必然带来互动，若画面质量和表达吸引力不足，AI 反而可能成为减分项。该案例从反面说明，真正关键的并不是 AI 使用得多不多，而是最终呈现质量是否过关。

上海某东方明珠 AI 重构视频总互动量为 18。该视频内容具有一定信息价值，但互动表现较低。该案例说明，信息实用价值未必能直接转化为短视频平台互动，在泛娱乐平台中，用户更容易被视觉冲击、趣味表达和情绪唤起吸引。这与回归中内容导向影响被风格、质感和节奏等更直接因素吸收的结果一致。

综合来看，典型案例与统计结果之间形成了较好的相互印证关系。高互动样本往往同时具备较高的 AI 技术质感、较鲜明的超现实表达和较紧凑的时长结构；低互动样本则更容易出现质感粗糙、表达平直或节奏拖沓的问题。因此，案例分析进一步强化了前文的核心判断，即 AI 文旅短视频的有效传播并不是简单依赖技术标签，而是依赖高质量视觉呈现、鲜明风格、较短时长和平台化表达共同作用。

同样带有 AI 标签的视频出现明显传播分化，主要与内容组合方式有关。高互动案例通常具备鲜明的地方符号、较顺滑的画面质感、与旅游场景相连的超现实表达、开头几秒内形成的视觉记忆点，以及能够引发转发和评论的话题性；低互动案例则更容易出现地方符号弱、画面质感粗糙、表达与旅游场景脱节或节奏拖沓等问题。

\section{实践建议}

地方文旅账号应避免把 AI 作为单纯的营销标签。本文结果表明，AI 技术质感才是影响传播效果的核心因素。实践中，应优先保证画面清晰度、人物和景观合理性、色彩协调性、运动自然度和文化符号准确性。对于生成痕迹明显、画面违和或内容粗糙的视频，应在发布前进行二次筛选和人工修正。

AI 的优势在于能够生成传统拍摄难以实现的场景。文旅账号可以围绕地方地标、历史文化、自然景观和城市符号设计抽象/超现实表达，形成具有平台传播力的视觉记忆点。将地方文化意象与城市景观结合，并通过 AI 形成奇观化画面，是一种可行的表达方向。但这类表达应与目的地特色相关，避免脱离地方文旅内涵。

视频时长显著负向影响传播效果，因此此类视频应尽量压缩无效铺垫，在开头快速呈现核心视觉亮点。对于以视觉奇观为主的视频，可以采用短时长、高密度、强节奏的表达方式；对于信息实用型内容，也应减少冗长解说，用更清晰的结构提升观看效率。

文旅传播不同于一般娱乐内容，目的地真实感和文化可信度仍然重要。AI 生成内容应服务于地方形象和旅游体验，而不是完全替代真实场景。实践中可以采用“真实素材+AI增强”“地方文化符号+AI再创造”的方式，在新奇感和真实感之间保持平衡。相关研究也表明，AI 工具的采用需要与文化保存、服务能力和受众信任共同考虑\cite{lit25}。

文旅账号应建立持续复盘机制，不只统计总互动量，还应区分点赞、评论、收藏和转发的不同含义。点赞代表即时认可，收藏代表信息保存价值，转发代表扩散潜力，评论代表讨论热度。运营者可以结合视频发布时间，观察单位时间互动表现，避免简单用累计互动量评价不同发布时间的视频。

进一步说，复盘机制应尽量从“事后看总量”转向“分阶段看过程”。例如，在内容发布后的前 24 小时、72 小时和一周内分别观察互动变化，可以更清楚地区分是首发阶段就表现强势，还是后期因转发和讨论逐步放大；对同类选题内容进行横向比较，则有助于判断问题究竟出在选题、画面质感、风格表达还是发布时间安排。对于官方文旅账号而言，这种结构化复盘既有助于提高 AI 内容生产效率，也有助于在有限资源条件下逐步积累更稳定的传播经验。
